\section{Ubicación del curso y tesis central}

Esta clase está impartida por el profesor Yu Su, de Ohio State University, y su tema son tres capacidades clave de los agentes de lenguaje (Language Agents): Reasoning, Memory y Planning. La clase no trata al Agent como una forma de producto única, sino como un objeto de sistema escalable: el modelo necesita percibir, decidir, ejecutar, revisar y mejorar de manera continua en un entorno dinámico.

\begin{importantbox}{Tesis central de esta clase}
\textit{``Reasoning is for better acting.''} El razonamiento no existe para exhibir pasos en un problema de papel, sino para elegir mejores acciones en un entorno real. En correspondencia, Memory no es una acumulación de bitácoras, y Planning tampoco es una lista estática de pasos; ambos están al servicio de la tasa de éxito y la estabilidad de las tareas a largo plazo.
\end{importantbox}

Yu Su comienza revisando los sesgos narrativos comunes en el auge de los Agent de los últimos dos años: por un lado, la expectativa externa de que ``2025 es el año uno de los Agent'' es sumamente alta; por otro, los equipos, en la práctica de producción, enfrentan con frecuencia problemas como costos fuera de control, colapsos en tareas de cadena larga e inconsistencia en los criterios de evaluación. El valor de este curso está precisamente en sustituir la narrativa de eslóganes por un marco de ingeniería.

\begin{warningbox}{Alta popularidad no equivale a alta disponibilidad}
\begin{itemize}
    \item Que una demo tenga éxito no significa un éxito estable bajo la distribución real de tareas;
    \item Que la calidad de una sola inferencia sea alta no significa que la calidad de la ejecución en múltiples rondas también lo sea;
    \item Que un modelo tenga puntuaciones más altas no significa que el costo total del sistema sea aceptable.
\end{itemize}
\end{warningbox}

\subsection{Resumen de la sección}
Esta sección deja establecido el criterio de juicio para toda la clase: la investigación de Agent debe centrarse en la calidad del comportamiento del sistema, y no en las capacidades de un modelo aislado. Las discusiones posteriores sobre Reasoning, Memory y Planning giran todas en torno a ``cómo elevar la tasa de éxito en tareas reales''.
